\section{Study Overview and Methods}

\subsection{Current evidentiary basis}

This version of the methods section is grounded in the materials currently archived in the repository: the site framing in \texttt{index.html}, the placement rationale in \texttt{bigredbutton\_placement.tex}, the corresponding bibliography, and study media under \texttt{studies/first-big-red-button-study/}. Those materials are sufficient to document the intended apparatus, procedure, and measures. They are not sufficient to claim that data collection is complete, because the repository snapshot does not contain participant logs, questionnaire exports, raw physiological traces, or finalized analysis scripts. The present section therefore documents the protocol truthfully and leaves all empirical outcomes to a later revision once those materials are available.

\subsection{Environment and apparatus}

The study situates participants in a neutral grey VR room containing a single large red button. Existing project instructions describe the object as view-locked, positioned on the horizontal midline and approximately 15 degrees below eye level, placed at roughly 0.75\,m from the viewpoint, and rendered at approximately 10\,cm diameter. These choices are consistent with work on comfortable object placement in head-worn displays and with anthropometric estimates of reachable arm's-length interaction \citep{Lin2021,Mirbagheri2024,Bonin2022,Alghofaili2019,McNamara2019}. The object's saturated red appearance against a low-contrast background further aligns with feature-singleton salience research \citep{Li2008,CastellottiEtAl2022}.

Repository instructions identify the headset as a Meta Quest 3 and the physiological sensor as a Polar H10 chest strap connected through Bluetooth Low Energy. One experimental condition uses live electrocardiographic timing to modulate a flashing red pulse or tone associated with the button. A control condition uses simulated ECG activity generated through NeuroKit2-style signal synthesis \citep{MakowskiEtAl2021}. Respiration derived from the same physiological stream is used as a secondary modulation channel for ambient room lighting, with inhalation brightening the environment and exhalation dimming it. The intended logic is to test whether the button can become more body-adjacent when both the object and the surrounding room respond to the participant's physiology \citep{RockstrohEtAl2019,LueddeckeFelnhofer2022,ChittaroEtAl2024}.

\subsection{Procedure and condition logic}

The study is framed as an agency-preserving encounter rather than a performance task. Participants enter the virtual room, hear a reassuring instruction track, and are told that they may press the button or refrain from pressing it at any time. This matters methodologically because it keeps non-pressing available as a valid outcome rather than treating abstention as failure. The project also includes a visible counter above the button that displays cumulative presses, making each action explicitly legible as data production.

The repository describes the live-ECG and sham-ECG conditions as double blind. That claim should be retained as a protocol intention rather than a verified implementation fact until the randomization and blinding procedures are archived in the repository. Likewise, the study materials indicate a two-session structure in which novelty and order are expected to matter. Participants are reportedly asked to re-encounter the button in a second round while acknowledging that the phenomenology of first contact cannot literally be recreated.

\subsection{Outcome families}

The current materials support four outcome families.

\begin{enumerate}
\item \textbf{Behavioral pressing metrics.} These include press count, time to first press, non-press as an outcome, and the temporal structure of inter-press intervals.
\item \textbf{Presence and spatial proximity.} The project records presence and felt peripersonal relation to the button, which is consistent with VR presence research and with multisensory accounts of peripersonal space \citep{CummingsBailenson2016,Serino2019,SerinoEtAl2018}.
\item \textbf{Subjective self-world relation.} The materials reference a measure of oneness inspired by psychedelic research and a custom companion measure of ``secondness'' intended to capture repetition and order effects. The oneness reference is therefore grounded in the revised Mystical Experience Questionnaire literature \citep{BarrettEtAl2015}, whereas secondness should currently be treated as a project-specific, not yet independently validated construct.
\item \textbf{Physiological and temporal-complexity analyses.} The site framing proposes entropy, recurrence, scaling, and burst analyses for press time series. Those analyses remain prospective until raw event logs and signal files are archived.
\end{enumerate}

\subsection{Analytic stance}

The key analytic problem is not merely whether people press more in one condition than another. The stronger question is whether physiological coupling changes the \emph{kind} of relation participants report and enact toward the button. Three comparison families follow from the current design logic: live versus sham physiological coupling, first encounter versus repeated encounter, and pressing versus non-pressing participants. Novelty is theoretically central here, because the first encounter with a pressable object may structure later experience more strongly than condition assignment itself \citep{KiddHayden2015,RanganathRainer2003}. Any later statistical analysis should therefore model order effects explicitly rather than treating them as nuisance variance.
